% Wersja źródłowa: audyt III, jednostka 01
\section{Przestrzenie liniowo-topologiczne: definicje i pierwsze przykłady}\label{sec:tvs}

\subsection{Cel i konwencje}

W analizie funkcjonalnej bada się przestrzenie liniowe, w których można
wykonywać przejścia graniczne, oraz odwzorowania liniowe zgodne z tymi
przejściami. Norma jest tylko jednym ze sposobów wprowadzenia topologii.
Na przykład topologia zbieżności punktowej funkcji czy naturalna topologia
funkcji gładkich nie muszą pochodzić od pojedynczej normy. Dlatego zaczniemy
od aksjomatów, które wiążą strukturę liniową z topologiczną.

W całym skrypcie $\K$ oznacza $\R$ albo $\C$ z ich zwykłą topologią.
Przestrzenie liniowe rozpatrujemy nad $\K$. Przez otoczenie punktu $x$
rozumiemy zbiór zawierający pewien zbiór otwarty zawierający $x$;
w szczególności otoczenie nie musi być otwarte. Przez $\mathcal N_X(x)$
oznaczamy rodzinę wszystkich otoczeń $x$ w przestrzeni topologicznej $X$.
Słowo „ciągły” zawsze odnosi się do rozważanych topologii, nie zaś
do niejawnie wybranej normy.

\begin{definicja}[Przestrzeń liniowo-topologiczna]\label{def:tvs}
Przestrzenią liniowo-topologiczną nad $\K$ nazywamy przestrzeń liniową
$X$ z topologią $\tau$, dla której odwzorowania
\[
 +\colon X\times X\longrightarrow X,\qquad (x,y)\longmapsto x+y,
 \qquad
 m\colon\K\times X\longrightarrow X,\qquad (\lambda,x)\longmapsto\lambda x
\]
są ciągłe, gdy w dziedzinach przyjmiemy topologie produktowe.
Nie wymagamy na razie warunku Hausdorffa.
\end{definicja}

Warunek ciągłości mnożenia przez skalar jest \emph{łączny} względem
$\lambda$ i $x$. Sama ciągłość każdej funkcji $x\mapsto\lambda x$,
dla ustalonego $\lambda$, jest warunkiem słabszym.

\begin{przyklad}[Topologia trywialna]
Niech $X$ będzie dowolną niezerową przestrzenią liniową.
Topologia $\{\varnothing,X\}$ czyni $X$ przestrzenią
liniowo-topologiczną: każde odwzorowanie o wartościach w $X$ jest
ciągłe. Nie jest ona jednak Hausdorffa. Przykład wyjaśnia,
dlaczego rozdzielamy aksjomaty przestrzeni liniowo-topologicznej
od aksjomatu rozdzielania punktów.
\end{przyklad}

\subsection{Operacje liniowe jako homeomorfizmy}

\begin{stwierdzenie}\label{prop:translation}
Niech $X$ będzie przestrzenią liniowo-topologiczną nad $\K$.
Dla każdego $a\in X$ przesunięcie $T_a(x)=x+a$ jest
homeomorfizmem $X$ na $X$. Dla każdego $\lambda\in\K\setminus\{0\}$
odwzorowanie $M_\lambda(x)=\lambda x$ jest homeomorfizmem $X$ na $X$.
\end{stwierdzenie}
\begin{proof}
Ciągłość $T_a$ wynika ze złożenia ciągłego odwzorowania
$x\mapsto(x,a)$ z dodawaniem. Jego odwrotnością jest ciągłe $T_{-a}$.
Analogicznie $M_\lambda$ jest ciągłe jako funkcja jednej zmiennej
otrzymana z mnożenia przez skalar, a jego odwrotnością jest
$M_{1/\lambda}$.
\end{proof}

\begin{wniosek}\label{cor:zero-neighborhood}
Dla dowolnej przestrzeni liniowo-topologicznej $X$, punktu $x\in X$
i podzbioru $U\subseteq X$ zachodzi równoważność:
$U\in\mathcal N_X(x)$ wtedy i tylko wtedy, gdy
$U-x\in\mathcal N_X(0)$. W szczególności topologia przestrzeni
liniowo-topologicznej jest całkowicie określona przez otoczenia zera.
\end{wniosek}
\begin{proof}
Korzystamy z tego, że $T_{-x}$ jest homeomorfizmem i
$T_{-x}(x)=0$.
\end{proof}

\begin{uwaga}
Jednorodność topologiczna nie oznacza, że każda przestrzeń
liniowo-topologiczna jest metryzowalna. Problem metryzowalności
będziemy rozstrzygać za pomocą baz otoczeń zera.
\end{uwaga}

\subsection{Otoczenia zera, zbiory zrównoważone i pochłaniające}

\begin{definicja}\label{def:balanced-absorbing}
Niech $X$ będzie przestrzenią liniową nad $\K$, a $A\subseteq X$.
Mówimy, że:
\begin{enumerate}[label=\textup{(\roman*)}]
 \item $A$ jest \emph{zrównoważony}, jeżeli $\lambda A\subseteq A$
       dla każdego $\lambda\in\K$ z $|\lambda|\leq1$;
 \item $A$ jest \emph{pochłaniający}, jeżeli dla każdego $x\in X$
       istnieje $t>0$ takie, że $x\in tA$;
 \item $A$ jest \emph{wypukły}, jeżeli $sx+(1-s)y\in A$
       dla wszystkich $x,y\in A$ i $s\in[0,1]$.
\end{enumerate}
W drugim punkcie $tA=\{ta:a\in A\}$, gdzie $t$ jest rzeczywistą
liczbą dodatnią, również gdy $\K=\C$.
\end{definicja}

\begin{uwaga}
Wszystkie trzy warunki są algebraiczne. Zrównoważenie nie
implikuje wypukłości. W analizie lokalnie wypukłej będziemy
konstruować bazy otoczeń zera złożone ze zbiorów jednocześnie
wypukłych i zrównoważonych.
\end{uwaga}

\begin{twierdzenie}[Podstawowe własności otoczeń zera]\label{thm:zero-neighborhoods}
Niech $X$ będzie przestrzenią liniowo-topologiczną nad $\K$.
Dla każdego $U\in\mathcal N_X(0)$ zachodzą następujące własności:
\begin{enumerate}[label=\textup{(\roman*)}]
 \item istnieje $V\in\mathcal N_X(0)$ takie, że $V+V\subseteq U$;
 \item istnieje otwarte, symetryczne otoczenie $V$ zera
       ($V=-V$) zawarte w $U$;
 \item istnieje otwarte, zrównoważone otoczenie $W$ zera
       zawarte w $U$;
 \item $U$ jest zbiorem pochłaniającym.
\end{enumerate}
\end{twierdzenie}
\begin{proof}
Wybierzmy otwarte otoczenie $U_0$ zera z $U_0\subseteq U$.
Ciągłość dodawania w punkcie $(0,0)$ daje otoczenia otwarte
$V_1,V_2$ zera takie, że $V_1+V_2\subseteq U_0$.
Dla $V=V_1\cap V_2$ otrzymujemy (i). Zbiór
$U_0\cap(-U_0)$ jest otwartym, symetrycznym otoczeniem zera;
to dowodzi (ii).

Dla (iii) oznaczmy przez $D=\{\lambda\in\K:|\lambda|\leq1\}$
domknięty dysk jednostkowy (dla $\K=\R$ przedział $[-1,1]$).
Zbiór $D$ jest zwarty. Dla każdego $\lambda\in D$, z ciągłości
mnożenia przez skalar w $(\lambda,0)$, istnieją otoczenia otwarte
$O_\lambda\subseteq\K$ punktu $\lambda$ oraz $V_\lambda\subseteq X$
zera takie, że $O_\lambda V_\lambda\subseteq U_0$.
Wybieramy skończone podpokrycie
$D\subseteq O_{\lambda_1}\cup\dots\cup O_{\lambda_m}$ i kładziemy
$V=\bigcap_{j=1}^m V_{\lambda_j}$. Wtedy $DV\subseteq U_0$.
Zbiór $A=DV$ jest zrównoważony, ponieważ $\mu D\subseteq D$
dla $|\mu|\leq1$. Ponieważ $1\in D$, zachodzi $V\subseteq A$,
a zatem $A$ jest otoczeniem zera. Wreszcie $W=\operatorname{int}A$
jest otwarte i zawiera zero. Jeśli $0<|\mu|\leq1$, to
$\mu W$ jest otwartym podzbiorem $\mu A\subseteq A$, wobec czego
$\mu W\subseteq\operatorname{int}A=W$. Dla $\mu=0$ mamy
$\mu W=\{0\}\subseteq W$. Tak więc $W$ jest zrównoważone.

Dla (iv) ustalmy $x\in X$. Funkcja
$\K\ni\lambda\mapsto\lambda x\in X$ jest ciągła i przyjmuje
wartość zero dla $\lambda=0$. Istnieje zatem rzeczywiste
$\varepsilon>0$ takie, że $\varepsilon x\in U$.
Wtedy $x\in\varepsilon^{-1}U$, czyli $U$ jest pochłaniający.
\end{proof}

\begin{przyklad}[Zrównoważenie bez wypukłości]
Zbiór $A=\{(x,y)\in\R^2:|x|\leq1,\ |y|\leq1\}\cup
\{(x,y)\in\R^2:|x|\leq2,\ |y|\leq\frac12\}$
jest zrównoważony, lecz nie jest wypukły:
$(1,1),(2,\tfrac12)\in A$, podczas gdy
$(\tfrac32,\tfrac34)\notin A$.
\end{przyklad}

\subsection{Warunek Hausdorffa i domknięcie zera}

\begin{twierdzenie}\label{thm:hausdorff-zero}
Niech $X$ będzie przestrzenią liniowo-topologiczną nad $\K$.
Następujące warunki są równoważne:
\begin{enumerate}[label=\textup{(\roman*)}]
 \item $X$ jest przestrzenią Hausdorffa;
 \item zbiór $\{0\}$ jest domknięty w $X$;
 \item $\bigcap_{U\in\mathcal N_X(0)}U=\{0\}$.
\end{enumerate}
Ponadto $N=\cl{\{0\}}$ jest domkniętą podprzestrzenią liniową
$X$ oraz
\[
 N=\bigcap_{U\in\mathcal N_X(0)}U.
\]
\end{twierdzenie}
\begin{proof}
Implikacja (i)$\Rightarrow$(ii) zachodzi w każdej przestrzeni
Hausdorffa. Aby wykazać wzór na $N$, ustalmy $x\in X$.
Jeżeli $x\in N$, to każde otoczenie $x$ zawiera zero.
Dla dowolnego symetrycznego otoczenia $V$ zera oznacza to,
że $0\in x+V$, czyli $-x\in V$, a z symetrii także $x\in V$.
Ponieważ symetryczne otoczenia tworzą bazę, $x$ należy do
wszystkich otoczeń zera.

Odwrotnie, jeśli $x$ należy do wszystkich otoczeń zera,
to dla każdego otoczenia $x+V$ punktu $x$, gdzie $V$ jest
symetrycznym otoczeniem zera, mamy $-x\in V$, skąd
$0\in x+V$. Oznacza to $x\in N$. Udowodniliśmy wzór.

Jeżeli (ii), to $N=\{0\}$ i powyższy wzór daje (iii).
Jeżeli (iii), to dla dowolnych $x\ne y$ istnieje otoczenie
$U$ zera, do którego nie należy $y-x$. Dobierzmy symetryczne
otoczenie $V$ zera takie, że $V+V\subseteq U$. Otoczenia
$x+V$ i $y+V$ są rozłączne: w przeciwnym razie
$y-x\in V-V=V+V\subseteq U$. To daje (i).

Pozostaje wykazać, że $N$ jest podprzestrzenią.
Jeśli $x,y\in N$ i $U$ jest otoczeniem zera, wybierzmy
otoczenie $V$ z $V+V\subseteq U$. Wtedy $x,y\in V$,
a więc $x+y\in U$. Ponieważ $U$ było dowolne, $x+y\in N$.
Dla $\lambda\ne0$ i $x\in N$, zbiór $\lambda^{-1}U$
jest otoczeniem zera, więc $x\in\lambda^{-1}U$,
czyli $\lambda x\in U$; dla $\lambda=0$ jest to oczywiste.
Zatem $N$ jest podprzestrzenią. Jest domknięta z definicji.
\end{proof}

\begin{uwaga}
Od tego miejsca \emph{przestrzeń liniowo-topologiczna} będzie,
o ile nie zaznaczymy inaczej, oznaczać przestrzeń liniowo-topologiczną
Hausdorffa. Warto jednak pamiętać, że wymaganie Hausdorffa jest
założeniem, a nie konsekwencją ciągłości działań liniowych.
\end{uwaga}

\subsection{Ciągłe odwzorowania liniowe}

\begin{twierdzenie}[Kryterium ciągłości w zerze]\label{thm:linear-continuity-zero}
Niech $X,Y$ będą przestrzeniami liniowo-topologicznymi nad tym samym
ciałem $\K$ i niech $T\colon X\to Y$ będzie odwzorowaniem liniowym.
Następujące warunki są równoważne:
\begin{enumerate}[label=\textup{(\roman*)}]
 \item $T$ jest ciągłe na $X$;
 \item $T$ jest ciągłe w punkcie $0$;
 \item dla każdego $V\in\mathcal N_Y(0)$ istnieje
       $U\in\mathcal N_X(0)$ takie, że $T(U)\subseteq V$.
\end{enumerate}
\end{twierdzenie}
\begin{proof}
Implikacja (i)$\Rightarrow$(ii) jest oczywista.
Warunek (iii) jest dokładnym zapisem definicji ciągłości w zerze,
ponieważ $T(0)=0$.
Jeśli (iii) oraz $x\in X$, to dla dowolnego otoczenia
$T(x)+V$ punktu $T(x)$ istnieje otoczenie $U$ zera takie,
że $T(U)\subseteq V$. Z liniowości
$T(x+U)=T(x)+T(U)\subseteq T(x)+V$.
Oznacza to ciągłość $T$ w punkcie $x$, a więc (i).
\end{proof}

\begin{wniosek}\label{cor:kernel-closed}
Niech $X$ będzie przestrzenią liniowo-topologiczną, $Y$
przestrzenią liniowo-topologiczną Hausdorffa, a $T\colon X\to Y$
ciągłym odwzorowaniem liniowym. Wtedy $\ker T$ jest domkniętą
podprzestrzenią liniową $X$.
\end{wniosek}
\begin{proof}
Jest podprzestrzenią z liniowości $T$. Ponieważ $Y$ jest
Hausdorffa, zbiór $\{0\}$ jest domknięty na mocy
twierdzenia~\ref{thm:hausdorff-zero}. Mamy
$\ker T=T^{-1}(\{0\})$, zatem jądro jest domknięte.
\end{proof}

\begin{przyklad}[Ciągła bijekcja liniowa bez ciągłej odwrotności]
Na $X=C([0,1],\R)$ rozważmy dwie normy:
$\|f\|_\infty=\max_{x\in[0,1]}|f(x)|$ oraz
$\|f\|_1=\int_{[0,1]}|f(x)|\,\mathrm d\lambda(x)$,
gdzie $\lambda$ jest miarą Lebesgue'a.
Identyczność $I\colon(X,\|\cdot\|_\infty)\to(X,\|\cdot\|_1)$
jest liniową ciągłą bijekcją, bo $\|f\|_1\leq\|f\|_\infty$.
Odwrotność nie jest ciągła. Istotnie, funkcje
$f_n(x)=\max\{1-nx,0\}$ spełniają $\|f_n\|_\infty=1$
i $\|f_n\|_1=1/(2n)$.
Nie przeczy to twierdzeniu o odwzorowaniu otwartym:
przestrzeń $(X,\|\cdot\|_1)$ nie jest zupełna. Rzeczywiście,
można dobrać ciąg funkcji ciągłych $g_n$ przyjmujących wartości
w $[0,1]$, równych zeru na $[0,\tfrac12-\tfrac1n]$ oraz jeden
na $[\tfrac12+\tfrac1n,1]$, liniowych pomiędzy tymi przedziałami
(dla $n\geq3$). Wtedy
$\|g_n-\chi_{[1/2,1]}\|_1\leq2/n\to0$.
W szczególności $(g_n)$ jest ciągiem Cauchy'ego w normie
$\|\cdot\|_1$. Gdyby zbiegał do $g\in C([0,1])$,
jednoznaczność granicy w $L^1$ dawałaby
$g=\chi_{[1/2,1]}$ prawie wszędzie. Z ciągłości $g$
wynikałoby wtedy $g=0$ na $[0,1/2)$ oraz $g=1$ na
$(1/2,1]$, co jest sprzeczne z ciągłością w $1/2$.
Zatem $(X,\|\cdot\|_1)$ nie jest zupełna.
\end{przyklad}

\subsection{Zbiory ograniczone bez użycia normy}

\begin{definicja}[Ograniczoność w sensie von Neumanna]\label{def:tvs-bounded}
Niech $X$ będzie przestrzenią liniowo-topologiczną nad $\K$.
Zbiór $B\subseteq X$ nazywamy \emph{ograniczonym}, jeżeli
dla każdego otoczenia $U$ zera istnieje $t>0$ takie,
że $B\subseteq tU$. W definicji nie wymaga się, aby $B$
był domknięty ani aby był otoczeniem zera.
\end{definicja}

\begin{lemat}\label{lem:bounded-test}
Niech $X$ będzie przestrzenią liniowo-topologiczną, a $B\subseteq X$.
Zbiór $B$ jest ograniczony wtedy i tylko wtedy, gdy dla każdego
zrównoważonego otoczenia $U$ zera istnieje $\varepsilon>0$ takie,
że $\lambda B\subseteq U$ dla wszystkich $\lambda\in\K$
z $|\lambda|\leq\varepsilon$.
\end{lemat}
\begin{proof}
Załóżmy, że $B$ jest ograniczony, i weźmy zrównoważone
otoczenie $U$ zera. Istnieje $t>0$ z $B\subseteq tU$.
Dla $\varepsilon=1/t$ i $|\lambda|\leq\varepsilon$
mamy $\lambda B\subseteq(\lambda t)U\subseteq U$.

Odwrotnie, niech $V$ będzie dowolnym otoczeniem zera.
Na mocy twierdzenia~\ref{thm:zero-neighborhoods} istnieje
zrównoważone otoczenie $U\subseteq V$. Z założenia
$\varepsilon B\subseteq U$ dla pewnego $\varepsilon>0$,
a zatem $B\subseteq\varepsilon^{-1}V$.
\end{proof}

\begin{twierdzenie}\label{thm:compact-bounded}
Niech $X$ będzie przestrzenią liniowo-topologiczną.
Wówczas:
\begin{enumerate}[label=\textup{(\roman*)}]
 \item każdy skończony podzbiór $X$ jest ograniczony;
 \item każdy zwarty podzbiór $X$ jest ograniczony;
 \item jeżeli $T\colon X\to Y$ jest ciągłym odwzorowaniem liniowym
       między przestrzeniami liniowo-topologicznymi, to obraz
       każdego ograniczonego zbioru $B\subseteq X$ jest ograniczony w $Y$.
\end{enumerate}
\end{twierdzenie}
\begin{proof}
Wystarczy udowodnić (ii) i (iii), ponieważ każdy zbiór
skończony jest zwarty. Ustalmy zwarty zbiór $K\subseteq X$
oraz otoczenie $U$ zera. Wybierzmy otwarte, zrównoważone
otoczenie $V\subseteq U$. Zbiór $V$ jest pochłaniający,
więc $X=\bigcup_{n\in\N,\,n\geq1}nV$: jeżeli $x\in tV$ dla
pewnego $t>0$, to po wybraniu liczby całkowitej $n\geq t$
mamy $x\in nV$, ponieważ $V$ jest zrównoważony.
Rodzina $(nV)_{n\geq1}$ jest wstępująca, a każdy $nV$
jest otwarty. Zwartość $K$ daje skończone podpokrycie,
a w konsekwencji pewne $n$ z $K\subseteq nV\subseteq nU$.
To dowodzi (ii).

Dla (iii) niech $W\in\mathcal N_Y(0)$.
Z ciągłości $T$ zbiór $T^{-1}(W)$ jest otoczeniem zera w $X$.
Istnieje $t>0$ z $B\subseteq tT^{-1}(W)$.
Z liniowości $T(B)\subseteq tW$, co kończy dowód.
\end{proof}

\begin{przyklad}[Ograniczoność normowa]
Niech $(X,\|\cdot\|)$ będzie przestrzenią unormowaną.
Zbiór $B\subseteq X$ jest ograniczony w sensie
definicji~\ref{def:tvs-bounded} wtedy i tylko wtedy, gdy
$\sup_{x\in B}\|x\|<\infty$ (przyjmujemy supremum zbioru
pustego równe zero). Istotnie, wystarczy sprawdzać
otoczenia postaci $\{x:\|x\|<r\}$, $r>0$.
\end{przyklad}

\subsection{Przykłady przestrzeni nienormowalnych}

\begin{przyklad}[Przestrzeń wszystkich ciągów]\label{ex:product-tvs}
Niech $X=\K^{\N}$ będzie przestrzenią wszystkich ciągów
$x=(x_1,x_2,\ldots)$, wyposażoną w topologię produktową.
Bazę otoczeń zera tworzą zbiory
\[
 U(F,\varepsilon)=\big\{x\in\K^{\N}:|x_j|<\varepsilon_j
          \text{ dla }j\in F\big\},
\]
gdzie $F\subseteq\N$ jest skończony, a $\varepsilon_j>0$ dla
$j\in F$. Dodawanie i mnożenie przez skalar są ciągłe,
ponieważ ich złożenia ze wszystkimi rzutami współrzędnymi
są ciągłe. Zatem $X$ jest przestrzenią liniowo-topologiczną Hausdorffa.

Nie jest ona normowalna: żadne otoczenie zera nie jest ograniczone.
Istotnie, niech $U(F,\varepsilon)$ będzie bazowym otoczeniem
zera i wybierzmy $j\notin F$. Weźmy otoczenie
$W=\{x:|x_j|<1\}$. Dla każdego $t>0$ ciąg
$x=(t+1)e_j$ należy do $U(F,\varepsilon)$, ale nie należy
do $tW$. Stąd $U(F,\varepsilon)$ nie jest ograniczone,
a więc żadne większe otoczenie zera nie jest ograniczone.
Gdyby topologia pochodziła od normy, otwarta kula jednostkowa
byłaby ograniczonym otoczeniem zera, co prowadzi do sprzeczności.

Warto zauważyć, że $e_n\to0$ w topologii produktowej:
dla każdego ustalonego $j$ współrzędna $(e_n)_j$
jest równa zero dla wszystkich dostatecznie dużych $n$.
\end{przyklad}

\begin{przyklad}[Topologia zbieżności punktowej funkcji]
Niech $S$ będzie zbiorem, a $X\subseteq\K^S$ podprzestrzenią
liniową, np. przestrzenią wszystkich funkcji ciągłych na
$S=[0,1]$. Topologia zbieżności punktowej na $X$
jest topologią podprzestrzeni odziedziczoną z produktu
$\K^S$. Ciąg (lub ogólniej sieć) funkcji zbiega w tej
topologii wtedy i tylko wtedy, gdy zbiega w każdym punkcie $S$.
W szczególności topologia ta nie wymaga normy supremum.
\end{przyklad}

\begin{przyklad}[Nieciągły funkcjonał na przestrzeni unormowanej]
Rozważmy $X=\R[t]$, przestrzeń wielomianów rzeczywistych
na $[0,1]$, z normą supremum. Funkcjonał liniowy
$L(p)=p'(0)$ nie jest ciągły. Dla $p_n(t)=(1-t)^n$
mamy $\|p_n\|_\infty=1$, ale $|L(p_n)|=n$.
Gdyby $L$ był ciągły w zerze, z liniowości otrzymalibyśmy
$|L(p)|\leq C\|p\|_\infty$ dla pewnego $C>0$ i wszystkich
$p\in X$, co przeczy powyższemu rachunkowi.
\end{przyklad}

\subsection{Przykład zasadniczy: \texorpdfstring{$L^p$}{Lp} dla \texorpdfstring{$0<p<1$}{0 < p < 1}}

W tym przykładzie zakładamy znajomość elementarnych pojęć miary
Lebesgue'a. Ma on pokazać, dlaczego lokalna wypukłość będzie
później odgrywać kluczową rolę.

\begin{twierdzenie}[Zupełna przestrzeń liniowo-topologiczna
bez niezerowych funkcjonałów ciągłych]\label{thm:lp-trivial-dual}
Niech $0<p<1$, a $X=L^p([0,1],\lambda)$ oznacza przestrzeń
klas funkcji mierzalnych $f\colon[0,1]\to\K$ takich, że
$\int_{[0,1]}|f(x)|^p\,\mathrm d\lambda(x)<\infty$,
gdzie utożsamiamy funkcje równe prawie wszędzie.
Wówczas wzór
\[
 d(f,g)=\int_{[0,1]}|f(x)-g(x)|^p\,\mathrm d\lambda(x)
\]
definiuje zupełną metrykę na $X$ zgodną ze strukturą
przestrzeni liniowo-topologicznej. Co więcej, każdy
ciągły funkcjonał $\K$-liniowy $\Lambda\colon X\to\K$
jest równy zeru.
\end{twierdzenie}
\begin{proof}
Dla liczb rzeczywistych $a,b\geq0$ i $0<p<1$ zachodzi
$(a+b)^p\leq a^p+b^p$. W połączeniu z nierównością
$|f-h|\leq|f-g|+|g-h|$ daje to nierówność trójkąta dla $d$.
Pozostałe aksjomaty metryki wynikają z własności całki.
Dodawanie jest ciągłe na mocy
$d(f_1+g_1,f_2+g_2)\leq d(f_1,f_2)+d(g_1,g_2)$.
Dla $\alpha,\beta\in\K$ oraz $f,g\in X$ mamy
\[
 d(\alpha f,\beta g)
 \leq |\alpha|^p d(f,g)+|\alpha-\beta|^p d(g,0).
\]
Jest to oszacowanie wystarczające do łącznej ciągłości
mnożenia przez skalar.

Wykażemy zupełność. Niech $(f_n)$ będzie ciągiem Cauchy'ego
w metryce $d$; wybierzmy mierzalnych reprezentantów jego wyrazów.
Możemy wybrać podciąg $(f_{n_k})$ tak, aby
$\sum_{k=1}^\infty d(f_{n_{k+1}},f_{n_k})<\infty$.
Z twierdzenia o monotonicznej zbieżności wynika, że
\[
 \sum_{k=1}^\infty|f_{n_{k+1}}(x)-f_{n_k}(x)|^p<\infty
 \quad\text{dla prawie każdego }x\in[0,1].
\]
Ponieważ $p<1$, zbieżność szeregu dodatnich liczb
$\sum_k a_k^p$ implikuje zbieżność $\sum_k a_k$
(wszystkie dostatecznie późne $a_k$ należą do $[0,1]$,
a wtedy $a_k\leq a_k^p$). Zatem szereg różnic
$\sum_k(f_{n_{k+1}}-f_{n_k})$ zbiega bezwzględnie
prawie wszędzie. Oznaczmy jego granicę powiększoną o
$f_{n_1}$ przez $f$. Z nierówności $(\sum a_k)^p\leq\sum a_k^p$
i lematu Fatou dostajemy
\[
 d(f_{n_k},f)
 \leq \sum_{j=k}^\infty d(f_{n_{j+1}},f_{n_j})\longrightarrow0.
\]
W szczególności $f\in L^p$ (bo $|f|^p\leq|f-f_{n_1}|^p+|f_{n_1}|^p$),
a ciąg Cauchy'ego $(f_n)$ także zbiega do $f$.

Pozostało zbadać przestrzeń ciągłych funkcjonałów liniowych. Niech
$\Lambda\colon X\to\K$
będzie ciągłym funkcjonałem liniowym. Ciągłość w zerze
daje $\delta>0$ takie, że
$d(f,0)<\delta^p$ implikuje $|\Lambda(f)|<1$.
Z jednorodności $d(af,0)=|a|^p d(f,0)$ wynika istnienie
$C>0$ takie, że
\[
 |\Lambda(f)|\leq C\Big(\int_{[0,1]}|f|^p\,\mathrm d\lambda\Big)^{1/p}
 \qquad(f\in X).
\]
Rzeczywiście, dla $f\ne0$ wystarczy zastosować poprzedni
warunek do $a f$ z $|a|=\delta/(2d(f,0)^{1/p})$.

Ustalmy dowolny zbiór mierzalny $A\subseteq[0,1]$.
Dla każdego $n\geq1$ można rozbić $A$ na rozłączne
zbiory mierzalne $A_1,\ldots,A_n$ o równej mierze
$\lambda(A_i)=\lambda(A)/n$; wynika to z bezatomowości
miary Lebesgue'a. Stąd
\[
 |\Lambda(\chi_A)|
 \leq\sum_{i=1}^n|\Lambda(\chi_{A_i})|
 \leq C\lambda(A)^{1/p}n^{1-1/p}\longrightarrow0,
\]
ponieważ $1-1/p<0$. Funkcjonał $\Lambda$ znika więc
na funkcjach charakterystycznych, a z liniowości również
na funkcjach prostych. Uzasadnijmy również używaną gęstość funkcji prostych.
Dla $f\in L^p([0,1],\lambda)$ obcięcie
$f_N=f\chi_{\{|f|\leq N\}}$ spełnia
$d(f_N,f)\to0$ na mocy zbieżności zmajoryzowanej.
Funkcja $f_N$ jest ograniczona; dzieląc kwadrat
$[-N,N]^2$ (dla $\K=\C$) albo przedział $[-N,N]$
(dla $\K=\R$) na dostatecznie małe komórki i zastępując
wartości $f_N$ reprezentantami komórek, otrzymujemy
mierzalne funkcje proste $s_k$ z
$|s_k-f_N|\leq1/k$ prawie wszędzie. Wtedy
$d(s_k,f_N)\leq k^{-p}\lambda([0,1])\to0$.
Funkcje proste są więc gęste w $L^p$; z ciągłości
$\Lambda=0$ na całym $X$.
\end{proof}

\begin{uwaga}\label{rem:local-convexity-purpose}
Na przestrzeni $L^p([0,1])$ z $0<p<1$ nie można oddzielić
niezerowego wektora od zera za pomocą ciągłego funkcjonału
liniowego: przestrzeń ciągłych funkcjonałów liniowych jest trywialna.
W dalszych
rozdziałach udowodnimy, że w każdej przestrzeni lokalnie
wypukłej Hausdorffa ciągłe funkcjonały liniowe oddzielają
punkty. Zatem powyższa przestrzeń \emph{nie jest lokalnie wypukła},
choć jest metryzowalna i zupełna.
\end{uwaga}

\subsection{Co zostało ustalone, a co wymaga dodatkowych założeń?}

W dowolnej przestrzeni liniowo-topologicznej możemy mówić o
ciągłości operatorów liniowych, zbiorach ograniczonych, otoczeniach
zera, podprzestrzeniach i zbieżności. Nie potrzebowaliśmy do tego
ani normy, ani iloczynu skalarnego, ani zupełności.

Nie możemy natomiast bez dalszych założeń zakładać, że istnieją
niezerowe ciągłe funkcjonały liniowe. Twierdzenie
\ref{thm:lp-trivial-dual} pokazuje, że jest to rzeczywista
przeszkoda. W kolejnych rozdziałach zbadamy podprzestrzenie, produkty
i ilorazy, następnie zbieżność sieci, ograniczoność oraz
metryzowalność. Później wprowadzimy seminormy i przestrzenie
lokalnie wypukłe.

\begin{uwaga}[Dalszy kierunek skryptu]
Pochodne Gâteaux i Frécheta nie wymagają teorii dystrybucji.
Wprowadzimy je w osobnym bloku dotyczącym przestrzeni Banacha
i Hilberta, a następnie wykorzystamy w twierdzeniu o funkcji
uwikłanej, zasadach wariacyjnych i rachunku na przestrzeniach
funkcyjnych. Dla przestrzeni zespolonych wyraźnie odróżnimy
różniczkowalność w sensie rzeczywistym od pochodnej
zespolono-liniowej.
\end{uwaga}
